\section{魏：最先跑起来的中原国家，最先感到疲惫}

从安邑向东望去，黄河与汾水之间的道路把魏推到列国往来的中央。三家灭智、周室册命之后，魏接过的不只是旧晋的一部分城邑和田土，也接过了一个不再容许卿族各自经营的中原。西面隔河是秦，东面通齐，南北又与韩、赵、楚相接；粮食、士人和兵车能够沿着这些道路汇集，敌军也同样能够沿着它们逼近。魏文侯一朝之所以在后世叙事中显得格外明亮，首先不在于某一位臣子忽然发明了新制度，而在于这个新国家较早地把田亩、法令、官吏和兵员视为同一场竞争里的资源。它以\eventref{WQ-0403-THREEJIN}{战国·三家分晋}所确认的名分为起点，却必须靠实际的征发与组织来守住这份名分。

\begingroup
\let\visibleeventanchor\eventanchor
\renewcommand{\eventanchor}[1]{\hypertarget{evt:#1}{}}
\renewcommand{\visibleeventanchor}[2]{\hyperlink{evt:#1}{#2}}

\begin{event}{前445年至前403年}{安邑、西河与新魏国的重心}{可靠纪年，早期政务细节有限}{魏}
\eventanchor{WQ-0445-WEI-WENHOU-ASCEND}\visibleeventanchor{WQ-0445-WEI-WENHOU-ASCEND}{战国·魏文侯即位}\eventanchor{WQ-0444-WEI-COUNSELORS}\eventanchor{WQ-0432-WEI-ANYI-URBAN-NETWORK}
魏文侯即位时，安邑并非一座只为朝会而存在的旧都。河东的农业、盐池资源以及通向黄河渡口和西河的转输路线，使它能够把近处的粮食、盐利和远处的戍守联结起来。这里也是魏与旧晋地理关系最紧的一面：它占据的是晋国腹地的一部分，不是坐拥天然边界的完整领土；向西越河要面对秦，向北、向东的扩张又会触动赵、中山和齐。故而“中原居中”既是可集中资源的条件，也意味着每一条道路都可能变成敌军的进路。

《史记·魏世家》与《资治通鉴》能确证文侯即位和魏氏经营西河的大致脊柱；安邑古城、盐池和交通研究则帮助解释这一选择的物质条件。它们不能给出一份魏国收入的总账，也不能把后来大梁的繁盛倒投给安邑。比较稳妥的判断是，魏在名分获得承认以前，已须将卿族的据点、河东的田地和边地的兵站编入同一套供给网络；周室的册命确认了既成的力量，却没有替它完成这件事。
\end{event}

\begin{event}{约前5世纪后期}{邺地的地方官与漳水记忆}{约年，传记叙事层次复杂}{魏}
\eventanchor{WQ-0435-WEI-YE-GOVERNANCE}\visibleeventanchor{WQ-0435-WEI-YE-GOVERNANCE}{战国·西门豹治邺}\eventanchor{WQ-0403-WEI-XIMENBAO-YE}
西门豹治邺的故事，留住的不是一份可以逐项核销的县令簿，而是一种地方治理的难题。魏要在赵魏交界的邺地得到稳定的田赋和兵役，就须处理漳水、乡里祭祀与地方摊派交错的秩序。《史记·滑稽列传》写河伯祭祀的整治，《吕氏春秋·乐成》也保存相关传说；后世因而把废除祭祀性勒索、兴修水利和整饬赋役归结为一位能吏的连贯举措。

这份记忆可以说明魏国的国家能力必须落在城邑与水边，而不能证明某年某月完成了哪一段渠工，更不能把传记中富于讽喻的情节改写成逐年行政实录。若说西门豹事迹有何历史位置，正在于它提示了地方官需要把原先由豪强、巫祝和乡里支配的成本重新纳入县邑秩序；至于其任官年代、河工次序和田赋数额，都应保留为未能坐实的部分。
\end{event}

\begin{event}{约前423年至前407年}{北向中山：胜利也要靠粮道等待}{年代范围较可靠，战事细节多有异说}{魏、中山}
\eventanchor{WQ-0423-WEI-LE-YANG}\visibleeventanchor{WQ-0423-WEI-LE-YANG}{战国·乐羊伐中山}\eventanchor{WQ-0421-WEI-ZHONGSHAN-SIEGE}\eventanchor{WQ-0408-ZHONGSHAN}\eventanchor{WQ-0407-WEI-ZHONGSHAN-ANNEX}
乐羊率军北攻中山，围城久而不下，最后使中山政权一度覆亡，是魏文侯时期最能显示远征成本的一组记录。它不是一支精兵到达便可结束的战役：山地城邑能够依托旧部众抵抗，魏军则须把兵员、粮食和将领的信任一并维持。《史记·魏世家》与《吕氏春秋·乐成》都保存了这场战争的轮廓，却对出兵年、围攻年次及乐羊故事的细节并不一致；著名的羹肉故事尤其不宜当作战场实录。

魏取得的只是阶段性北扩，而非没有后患的吞并。把主要城邑纳入控制，意味着还要安排道路、守备和供应；中山后来能够复起，赵魏之间的边邑摩擦也会加深。因而中山之役并非“魏武卒无敌”的预告，而是较早展示了战国扩张中攻城、转输和占领同样重要。
\end{event}

\eventanchor{WQ-0413-WEI-QIN-HEXI}
\begin{event}{前413年至前405年前后}{西河的渡口、道路与边防行政}{可靠纪年与后设推断并存}{魏、秦、郑}
\visibleeventanchor{WQ-0413-WEI-QIN-HEXI}{战国·魏争河西}\eventanchor{WQ-0412-WEI-ZHENG}\eventanchor{WQ-0410-WEI-HEXI-OFFICES}\eventanchor{WQ-0408-XIHE-PROCESS}\eventanchor{WQ-0405-WEI-HEXI}\eventanchor{WQ-0405-WEI-XIHE-ROAD}
魏军沿黄河向河西城邑推进，同时向郑地方向施压，测试通往洛阳东南的交通线；这两件事共同显示魏不是只守安邑。西河的意义也不只是地图上的一块边地：渡口要有人戍守，远征的粮食要沿河东道路送达，前线城邑还须有人传令、征发和维持秩序。后世所说的“西河经营”，所指应是这样的连续过程，而非某一次取城即告完成。

《史记·魏世家》和《资治通鉴》提供了河西竞争的大致年代；河西古道、渡口和历史地理研究补足空间关系。魏代守邑官名、道路工程和具体征粮程序却没有连续档案，不能拿睡虎地秦简所见的成熟秦制替魏逐项补齐。能够说的只是：西河若没有常设的军政网络，魏在中山与秦之间反复投射兵力便无从谈起；这种网络越向边地延伸，维护它的财政和人力成本也越高。
\end{event}

\eventanchor{WQ-0408-WEI}
\begin{event}{前408年前后}{李悝传统与魏国的早期整顿}{年代与细节有争议}{魏}\visibleeventanchor{WQ-0408-WEI}{战国·魏国早期改革}\eventanchor{WQ-0420-WEI-LAW-JUDGMENT}
《史记·魏世家》把李悝列在魏文侯所任用的人物之中；较晚的《汉书·食货志》则把“尽地力”与平籴的议论系于他。两种叙事共同保存了一项重要的历史记忆：当谷物丰歉足以摇动价格、徭役和军粮时，国家不能只在收成之后向封邑索取，而要设法看见土地的产出、储藏与流通。魏的优势正在这里逐步形成：旧贵族的门第仍有分量，但能被官府统计、调度和供给的资源，开始比家族私下的从属关系更能决定一国的行动速度。

不过，后人常把这段历史收束成“李悝作《法经》、魏国遂变法”的整齐开端，材料并不容许这样的肯定。今见《法经》篇目主要出自后出的法制史追述，原书早已不存；“平籴”的办法也不能据此倒推出一套已可逐条复原的魏国法规。从现存材料看，更稳妥的说法是：魏国早期确有把农政、刑名和任官纳入国家竞争的趋势，李悝成为这一趋势最有力的后世符号。个人的履历、政策的形成和制度的长期运行，不能被压成同一件事。
\end{event}

\begin{debate}[李悝与《法经》：传统可以指路，不能代替原件]
《史记》《汉书·食货志》提供了李悝与魏文侯、农政议论相连的叙述；《晋书·刑法志》等后出材料保存了《法经》六篇的法制传统。它们足以说明魏被记忆为早期制度竞争的实验场，却不足以确证每一条法文、每一次颁行和全部政策都出自李悝一人。出土秦律能够帮助理解先秦法律文书化的深度，但不是魏国改革的直接档案。
\end{debate}

李悝传统所指向的，是把散在乡里的收成与劳力纳入可计算秩序的努力；吴起的故事所指向的，则是把同一批可征发的人变成可反复使用的军队。两者未必是一个人设计的一套蓝图，却在魏的中原处境中彼此需要：没有较稳定的供给，精锐无从久驻；没有能在边境取胜的军队，田地与市场也难以免于他国征取。

\eventanchor{WQ-WEI-WUQI}
\begin{event}{魏文侯、武侯时期}{吴起与魏武卒}{年代与细节有争议}{魏}\visibleeventanchor{WQ-WEI-WUQI}{战国·魏武卒}\eventanchor{WQ-0400-WEI-WUZU}\eventanchor{WQ-0391-WEI-WUZU-DRILL}
《史记·孙子吴起列传》说吴起在魏用兵有功，后世遂以“魏武卒”概括魏军的锐利。《吴子》又描绘负甲、持兵、行军的选拔条件，仿佛一张可以逐项核验的征募清单。可这部兵书的成书层累和这些标准的适用年代都存在疑问，不能把它直接当作魏文侯、武侯朝的一份军籍簿。可以把握的，是魏确实在战国早期以更严密的选拔、训练、赏罚和供给来提高军队的可用性；至于“武卒”的人数、装备和全部考核办法，仍须留在传说与制度史之间审慎辨认。

这里的人事转折并不神秘。像吴起这样的外来行动者能够在魏获得位置，反映的是君主需要能够组织战争的专业才能；而这类才能一旦与旧有贵族、君权和军功分配发生摩擦，也不可能只凭战场声望安然存在。史书随后写吴起离魏入楚，所保留的正是这种政治处境的结果，而非一段可以补写出心理独白的个人传奇。
\end{event}

\eventanchor{WQ-0389-YINJIN}
\begin{event}{前389年及其后}{阴晋胜利：把战场成果留在渡口}{年代与战果数字有争议}{魏、秦}
\visibleeventanchor{WQ-0389-YINJIN}{战国·阴晋之战}\eventanchor{WQ-0389-WEI-QIN-YINJIN-AFTERMATH}
《史记·孙子吴起列传》把吴起在阴晋击秦保存为魏军早期强盛的证据，《吴子》又使后人更容易把它读成武卒制度的必然胜利。胜负的基本脊柱可以把握，战果数字、战场地望和兵种构成却仍有争议。尤其不能因为后出兵书描绘严密的考选，便断言每一名参战者都按同一套标准编入军籍。

阴晋更值得注意的后果，是魏若要把一次胜利变成对秦的压力，便须继续守住黄河渡口和粮道。史书没有保存胜后增设戍守的明文；将胜利与渡口控制联系起来，是根据边防竞争所作的有限推断，而非可复原的建置清单。正是这种从野战到守邑的转换，才显示早期魏军的能力不是单纯的冲锋技艺。
\end{event}

\eventanchor{WQ-0396-WEI-WENHOU}
\begin{event}{前396年至前362年前后}{继承裂缝与大梁的选择}{王年较可靠，迁都过程有异说}{魏、韩、赵}
\visibleeventanchor{WQ-0396-WEI-WENHOU}{战国·魏文侯卒}\eventanchor{WQ-0396-WEI-WENHOU-FUNERAL}\eventanchor{WQ-0395-WEI-WU-HOU-COURT}\eventanchor{WQ-0373-WEI-ZHAO-FRICTION}\eventanchor{WQ-0370-WEI-SUCCESSION}\eventanchor{WQ-0370-WEI-SUCCESSION-ESTATES}\eventanchor{WQ-0370-HAN-ZHAO-WEI-INTERVENE}\eventanchor{WQ-0369-WEI-HUIWANG}\eventanchor{WQ-0362-WEI-DALIANG-CAPITAL}\eventanchor{WQ-0349-WEI-DALIANG-WATERWAYS}
文侯死于前396年，武侯承继安邑朝廷并延续西河路线；《史记·魏世家》能够给出这条王位线，却不能使丧礼程序或旧臣名单忽然清楚。前370年武侯死后，继承冲突终于让这种不确定变成公开的政治风险。韩、赵出兵干预，魏罃压倒公中缓而即位。传世材料足以确定邻国介入和惠王继位，关于双方怎样争夺封邑、军队与宫廷支持的细节则各有异文，不宜据此编排一场完整的宫廷政变。

惠王将政治重心移向大梁，是在这种内外压力下改变资源重心的选择。安邑靠近河东和西河，亦离秦的进攻方向更近；大梁连接鸿沟及黄河支流，更便于汇集粮食、车马和中原的使者。迁都的确年、完成步骤和早期渠工都难以从《魏世家》与后出的《水经注》复原。可以看见的却是一个战略取舍：水网扩大了漕运与守城的能力，也把魏日后的防务更深地系在可被敌方利用的河沟上。
\end{event}

魏军的锋芒使它能够向外施压，也使中原诸国更早学会以联盟回应它。前354年，魏军围赵都邯郸，赵向齐求援。齐不必在赵城下同魏军硬碰，便可选择把战事引回魏所倚重的腹地；《史记·孙子吴起列传》所叙田忌、孙膑的行动，由此成为“围魏救赵”的经典模型。战场在桂陵，胜负所打断的却是魏以一场围城迫使赵屈服的盘算。

\eventanchor{WQ-0354-GUILIN}
\begin{event}{前354年起}{邯郸危局与桂陵受挫}{可靠纪年}{魏、赵、齐}\visibleeventanchor{WQ-0354-GUILIN}{战国·桂陵之战}\eventanchor{WQ-0353-WEI-HANDAN-SUPPLY}
桂陵的故事之所以传得最广，正在于它把复杂的战略选择压缩为一条迂回路线。但围邯郸、齐军出动与桂陵交锋在不同史书中的编年并非完全整齐，后出的兵家叙事又强化了孙膑运筹的戏剧性。能够确定的是，齐的介入迫使魏在赵地之外顾及本土和补给线；\eventref{WQ-0354-GUILIN}{战国·桂陵之战}的直接后果，不是魏立刻失去一切，而是它再也不能把中原的每一场战争当成两国之间的单线较量。
\end{event}

十余年后，魏攻韩，齐又以救援者的身份出兵。马陵一役在前342年的传统编年中成为魏军受创最深的节点。传世叙事写孙膑减灶诱敌、庞涓自尽，画面锐利，却主要依赖《史记》与《战国策》的后出铺陈；银雀山汉简所见《孙膑兵法》证明相关兵学传统并非全然无根，却不能逐句为这场战役作证。应当保留的，不是一则奇谋的细节，而是魏军在长期外战、强邻合纵与补给压力交织之下，已难维持早期那种由它选择战场的主动。

\eventanchor{WQ-0342-MALING}
\begin{event}{前342年}{马陵之战与魏的战略收缩}{可靠纪年}{魏、韩、齐}\visibleeventanchor{WQ-0342-MALING}{战国·马陵之战}\eventanchor{WQ-0342-HAN-QI-AID}\eventanchor{WQ-0342-MALING-PRISONERS}\eventanchor{WQ-0341-WEI-MALING-RECOVERY}
马陵之后，魏仍是大国，后来还与齐互称王；败战并不等于国家骤然坍塌。它改变的是力量的相对位置：齐证明自己能够越出东方干预中原，韩、赵也不再只是魏的附庸式邻国，西面的秦则在河西与关中的竞争中等待机会。魏的地理曾使它最先聚拢道路和资源，如今也使它必须在更多方向上分兵、结交和付出财政成本。战国的中心没有消失，只是不再由魏独占。

马陵之后魏重编残军、依靠大梁仓储维持防线的说法，说明败者仍有恢复军政的能力；重编的程序、仓储数量和太子申被俘、庞涓战死等具体情节，却分别受《史记》《战国策》与后来的兵家叙事塑造，不能合并成一份无误的战报。
\end{event}

\begin{sourcenote}[桂陵、马陵与魏国衰势]
两役的基本对抗及其前后政治影响，可据《史记·孙子吴起列传》《魏世家》与《战国策》互相参照；人物谋略、伤亡数字和具体战术则带有显著的后世叙事加工。现代研究通常把目光放在齐魏竞争、行军与后勤条件、孙膑故事的成形过程，而不把“减灶”一类细节当作无需检验的战场实录。
\end{sourcenote}

\begin{event}{前330年至前302年}{失去河西后，魏如何同秦谈判}{可靠纪年为主，外交条件多不完整}{魏、秦}
\eventanchor{WQ-0330-HEXI}
\visibleeventanchor{WQ-0330-HEXI}{战国·魏割河西}\eventanchor{WQ-0330-WEI-HEXI-HOUSEHOLDS}\eventanchor{WQ-0319-WEI-XIANG}\eventanchor{WQ-0319-WEI-DEFENSIVE-ALLIANCE}\eventanchor{WQ-0319-HUISHI-WEI}\eventanchor{WQ-0311-WEI-XIANG-DIPLOMACY}\eventanchor{WQ-0303-QIN-PUBAN}\eventanchor{WQ-0303-QIN-PUBAN-WEI-RIVER-FORD}\eventanchor{WQ-0302-QIN-WEI-YINGTING}\eventanchor{WQ-0302-WEI-CHAOQIN}
前330年割让河西，并不是地图上抹去一片颜色而已。魏原来用来牵制秦的田邑、渡口和编户，自此成为秦的资源；安邑朝廷的西部财源也随之收缩。编户如何改隶并没有连续账簿，不能把人口、田亩和赋额说成精确的数字，但《史记·魏世家》《秦本纪》的割地记录足以说明，魏的西河体系从此不再是可以反攻的纵深，而成了秦向东经营的起点。

惠王死于前319年，魏襄王继位后更依赖使聘、局部结盟和游士的辩说来填补马陵后的兵力缺口。惠施参与魏廷的年月、职掌及其名辩究竟怎样转化为外交决策，主要须由《庄子》《吕氏春秋》和传世叙事拼合，不可把他塑成掌握全部国策的一位相。前311年前后秦惠文王去世，魏可以借西方权力交接调整防备，却只得到有限喘息。

前303年秦取蒲阪，黄河渡口的联络受威胁；次年魏王至应亭朝秦，秦归还蒲阪。两件事在编年上可以互证，交换的仪式、城邑实际控制时段和屈从条件却没有完整文书。它们仍显示魏的选择已不同于阴晋之后：它能以朝见换取局部缓和，却难以凭一场胜仗重建河西边防。
\end{event}

\begin{event}{前296年至前266年}{河东城链的蚕食与大梁的收缩}{可靠纪年与地望争议并存}{魏、秦、韩、赵、燕}
\eventanchor{WQ-0296-WEI-XIANG-DEATH}
\visibleeventanchor{WQ-0296-WEI-XIANG-DEATH}{战国·魏襄王卒}\eventanchor{WQ-0295-WEI-AI}\eventanchor{WQ-0294-WEI-HEDONG-REFUGEES}\eventanchor{WQ-0292-BAIQI-YUAN}\eventanchor{WQ-0290-WEI-HEDONG-CESSION}\eventanchor{WQ-0289-SIMACUO-YUAN-HEYONG}\eventanchor{WQ-0283-QIN-DALIANG-THREAT}\eventanchor{WQ-0282-QIN-AN-DAILIANG-AXIS}\eventanchor{WQ-0276-BAIQI-WEI-TWO-CITIES}\eventanchor{WQ-0266-QIN-WEI-XINGQIU-HUAI}\eventanchor{WQ-0266-WEI-DAILIANG-PRESSURE}
魏襄王卒、哀王继位而又卒的王位更替，留下的材料远比战争简略。真正持续压迫魏的，是\eventref{WQ-0293-YIQUE}{战国·伊阙之战}以后秦军取得的东进条件。韩魏受创、河东城邑受压，居民逃散和向内地避难很可能随之发生；不过迁徙规模没有可信统计，不能把“逃散”写成一场人数确定的迁徙。对魏而言，前沿失去的不只是城墙，也包括可征发的劳力、粮仓与通往大梁的缓冲。

此后取垣而还、司马错攻垣与河雍、前290年魏割河东、前283年秦取安城而燕赵援魏，组成了一条比单场大战更难逆转的压力线。传世记载让人看见秦以城、桥和河上节点切断魏的防线，却不允许我们精确画出每一座城的占领期限和每一次救援的规模。前266年邢丘、怀失守后，魏不得不把大梁西北的防御更多压缩在水网与近郊城邑之间。秦军的推进因而不是直线冲到都城，而是一层层减少魏可以等待援军和调度粮食的空间。
\end{event}

\begin{event}{前276年至前273年}{安釐王、信陵君与白起压力下的军政分配}{王年较可靠，兵数与人事细节有争议}{魏、秦、韩、赵}
\eventanchor{WQ-0276-WEIANXI}
\visibleeventanchor{WQ-0276-WEIANXI}{战国·魏安釐王即位}\eventanchor{WQ-0276-WEI-ANXI-COURT}\eventanchor{WQ-0276-XINLING-FENGJUN}\eventanchor{WQ-0276-WEI-DALIANG-FLOODWORKS}\eventanchor{WQ-0275-QIN-WEI-BAIQI}\eventanchor{WQ-0275-BAIQI-WEI-PRISONERS}\eventanchor{WQ-0275-WEI-MANGMAO-DEFENSE}\eventanchor{WQ-0275-WEI-POST-DEFEAT-RECRUITMENT}\eventanchor{WQ-0273-WEI-MANGMAO-RESCUE}
魏昭王死后，安釐王即位。连续失城使王室既要倚重宗室与将领的动员，又不能轻易放出独立的军权；无忌受封信陵，正是在这种条件下逐渐积累宾客、封号与宗室声望。封邑的范围和封授程序不足以复原，安釐王朝早期的人事任命也不全，但《魏世家》与《魏公子列传》足以说明，王、相与公子的军政关系已成为魏抵抗秦军时不可回避的内部问题。

白起前275年大破魏军，传本常以“斩首十万”制造震慑。这一王年可以把握，数字却不应当成为人口统计；《白起王翦列传》所塑造的恐惧也是秦军威势叙事的一部分。魏将芒卯组织河东与大梁守备，战败后又须重征兵员、修补城防并向赵楚求援，都是合乎记录的压力反应，具体征发数额和援军承诺则不能实数化。前273年芒卯会赵援韩于华阳而未能阻止秦军各个击破，显示联军能把危险延后，却没有解决谁来统一指挥、谁来负担补给的问题。
\end{event}

\begin{event}{前258年至前257年}{邺地驻军、兵符与救赵}{可靠纪年与列传戏剧化并存}{魏、赵、秦}
\eventanchor{WQ-0258-WEI-REFUSES-AID}
\visibleeventanchor{WQ-0258-WEI-REFUSES-AID}{战国·魏廷犹疑救赵}\eventanchor{WQ-0258-WEI-JINBI-AT-YE}\eventanchor{WQ-0257-XINLING}\eventanchor{WQ-0257-XINLING-PRIVATE-COMMAND}\eventanchor{WQ-0257-ZHUHAI-KILLS-JINBI}
秦围邯郸时，魏若全力救赵，便担心秦军转而直压大梁；若坐视赵亡，又会失去中原东北的一道屏障。魏安釐王命晋鄙率军而令其驻邺观望，把这两种风险悬在同一道军令上。这个驻军选择在《史记·魏公子列传》与《资治通鉴》的编年中可见，其具体传令与兵符流转主要仍来自列传，不能被误当成一份保存完整的军令文书。

无忌依靠私门网络取得兵符、朱亥杀晋鄙、魏军转而救赵，是魏国最著名的一次非常调兵。它使赵都危局得以缓解，也使魏获得联军解围的结果；同时，它把常规王命和宗室私人动员之间的裂缝暴露得再清楚不过。侯嬴、朱亥的言行和兵符故事有高度文学化的传述，应当辨别其叙事加工；但非常规指挥权确实改变了这支魏军的去向，且为魏王后来猜忌信陵君提供了现实的政治理由。
\end{event}

\begin{sourcenote}[从河西到大梁：不能由城名补成精确疆界]
前330年以后秦魏关系的年代节点，主要可据《史记·魏世家》《秦本纪》及《资治通鉴》相互校读；伊阙、垣、河雍、安城、邢丘、怀和蒲阪的地望与控制期限，则仍须由历史地理研究不断核对。魏的编户转隶、难民数量、河工形制和城防配置没有完整魏国档案，正文所说的财源收缩、人口流动与防御内缩，都是从割地、取城和求援的连续记录作出的有限推论。
\end{sourcenote}

\begin{event}{前247年至前225年}{能合纵的人离去，水网都城被反用}{可靠纪年为主，围城工程细节不可复原}{魏、秦、诸侯}
\eventanchor{WQ-0253-XINLING-EXILE}
\visibleeventanchor{WQ-0253-XINLING-EXILE}{战国·信陵君滞赵}\eventanchor{WQ-0247-XINLING-RETURNS-WEI}\eventanchor{WQ-0244-QIN-WEI-APPROACH}\eventanchor{WQ-0244-MENGAO-WEI-CHANG}\eventanchor{WQ-0244-WEI-DALIANG-WATER-GATES}\eventanchor{WQ-0243-WEI-ANXI-DEATH}\eventanchor{WQ-0243-WEI-XINLING-DEATH}\eventanchor{WQ-0243-WEI-XINLING-FUNERAL}\eventanchor{QIN-0228-WEI-JIA}\eventanchor{WQ-0225-WEI-DALIANG-INUNDATION}\eventanchor{WQ-0225-WEI-JIA-SURRENDER}
救赵之后，信陵君长期留赵；他的滞赵时间与魏王态度被列传放大到什么程度仍可讨论，但私调大军触发猜忌这一政治处境不难理解。前247年秦军压迫河外，魏王又不得不召回无忌，以其宾客与声望联络诸侯。五国在\eventref{WQ-0241-HEZONG}{战国·五国合纵}中再度攻秦而未能突破函谷关，不能被写成一支受单一指挥的常备联军；它更像各国在秦已取得东郡、直接威胁中原后的一次困难协调。信陵君死于前243年，魏由私门维系的跨国联络再难集中到一人手中。

同年安釐王卒，此后王室更替并未带来新的外部纵深。蒙骜前244年取畅、有诡等城，秦军沿大梁西北的城链继续试压。魏在城壕、水门和仓储上作出整治，是面对陆上屏障缩短的合理反应；大梁水系的工程年次和形制却难精确归属，不能把后来所见的城防倒推为一次完成的惠王或安釐王工程。前228年魏王假继位时，魏仍有王号与都城，但守都已成为其最主要的战略。

前225年，王贲引河沟水围灌大梁。强攻水网都城的代价很高，利用水系正是秦军绕开这一代价的办法；筑堤、决水、城内伤亡和受淹范围则没有可独立核验的细账。魏王假请降，结束的是魏国的王室与都城防御体系，不是一种从来没有代价的“地理宿命”。大梁曾让魏集中漕运、市场和军队，到最后也使秦能把长期经营的水利条件变成围城工具。
\end{event}

\begin{event}{前209年至前208年}{魏名的短暂复起}{年代较可靠，拥立和围城细节有争议}{魏地、秦}
\eventanchor{QIN-0209-WEIJU}
\visibleeventanchor{QIN-0209-WEIJU}{秦末·魏咎称王}\eventanchor{QIN-0209-WEI-JIU-RESTORATION}\eventanchor{QIN-0208-WEI-JIU-SIEGE}\eventanchor{QIN-0208-ZHOUSHI-WEI-RETREAT}
魏亡并未使“魏”立刻从关东政治中消失。秦末起事中，魏咎被拥立为魏王，周市等人借旧都与旧宗室的名义组织魏地力量；这说明战国国家留下的不仅是城墙，也有地方可辨认的政治记忆。可是这种复国势力的军政资源远不能与文侯、惠王时相比。章邯围临济，魏咎死守而亡，周市所部又在秦军反攻下放弃部分城邑、转而支持魏豹。拥立地点、人物关系、死因和撤退路线都有异文，却不妨碍我们看见：旧国名号能够迅速聚拢支持，不能自动恢复已经失去的财源、城防与统一指挥。
\end{event}

\begin{turningpoint}[制度得失：先行者为何也会落后]
魏较早把农、法、兵连成国家能力，因此取得了先发的速度；承担代价的却是被征发的乡里、持续供给军队的财政，以及必须面对多面威胁的中原社会。桂陵与马陵不是制度失效的简单判决，而是地缘、联盟和战争消耗对先发优势的反馈。河西的割让、河东城链的收缩和大梁的水围，则让这种反馈落到税源、渡口、兵员与都城空间上。

李悝、《法经》与吴起、武卒之所以始终重要，是因为后出的法律和兵学传统记住了魏曾经提出的问题：怎样让土地、官吏、装备和战功为国家所用。它们不能替代魏国已经执行的法规、名册或工程记录。魏未能把早期的组织能力转化为不受制衡的霸权，也不意味着它什么也没有留下；秦所面对并重组的，正是这种将财政、行政和军队系在领土竞争上的难题。
\end{turningpoint}
\endgroup
